\section{Limitations of the Study}
%szb: sztem ez külön chapter szokott lenni

This study has several limitations that should be considered when interpreting
the results.

First, the experiments focused only on Java-to-C++ translation. Although this
language pair is representative because of the syntactic and semantic
differences between the two languages, the findings may not generalize to other
language pairs. Previous studies have shown that model performance can vary
substantially depending on the source and target languages, particularly when
the languages differ in programming paradigm or available training data
\cite{llm_limitations}.

Second, only two models were evaluated: StarCoder2 and
Qwen2.5-Coder. Although these are strong open-source code models, other models
such as CodeLlama, DeepSeek-Coder, or larger proprietary systems may exhibit
different behavior. Prior work suggests that larger models generally achieve
better functional correctness, but they also require substantially more
computational resources \cite{llm_survey}. This means that the observed results
may partly reflect the limitations of the selected model size rather than the
overall capability of large language models.

Another limitation is that the fine-tuning process relied exclusively on the
XLCoST dataset. While XLCoST provides aligned Java and C++ code examples, it
still covers only a subset of possible coding styles, programming domains, and
problem types. Dataset quality and diversity are known to have a strong impact
on model performance, especially for code-related tasks \cite{zhu2022xlcost}.

The evaluation process also has limitations. Only one execution-based benchmark,
CodeEditorBench, was used to assess functional correctness. Although
CodeEditorBench provides realistic evaluation scenarios, other benchmarks such
as HumanEval, MBPP, or PolyHumanEval may produce different results because they
focus on different types of tasks and evaluation criteria \cite{chen2021codex,
austin2021mbpp, mhumaneval}. %szb: de hisz eezek nem is codetanslation benchmarkok. A plyHE-t kivéve, azt nem ismerem, az lehet az, nem tudiom - de a másik 2 fix nem
Therefore, the conclusions of this thesis are limited to the
specific benchmark setup that was used.

Finally, hardware limitations may have influenced the quality of the
fine-tuning process. Due to memory and computational constraints, only 3B
parameter models were used and parameter-efficient fine-tuning techniques were
applied. Larger models and longer training schedules may have achieved better
results, but they would require significantly more computational resources
\cite{hu2021lora}.